
\documentclass{article}
\usepackage{colortbl}
\usepackage{makecell}
\usepackage{multirow}
\usepackage{supertabular}

\begin{document}

\newcounter{utterance}

\centering \large Interaction Transcript for game `dond', experiment `coop\_en', episode 17 with Qwen3{-}32B{-}t0.0.
\vspace{24pt}

{ \footnotesize  \setcounter{utterance}{1}
\setlength{\tabcolsep}{0pt}
\begin{supertabular}{c@{$\;$}|p{.15\linewidth}@{}p{.15\linewidth}p{.15\linewidth}p{.15\linewidth}p{.15\linewidth}p{.15\linewidth}}
    \# & $\;$A & \multicolumn{4}{c}{Game Master} & $\;\:$B\\
    \hline

    \theutterance \stepcounter{utterance}  
    & & \multicolumn{4}{p{0.6\linewidth}}{
        \cellcolor[rgb]{0.9,0.9,0.9}{
            \makecell[{{p{\linewidth}}}]{
                \texttt{\tiny{[P1$\langle$GM]}}
                \texttt{You are playing a cooperative negotiation game in which you have to agree on how to divide a set of items among you and another player.} \\
\\ 
\texttt{Rules:} \\
\texttt{(a) You and the other player are given a set of items. Each of you is also given a secret value function, representing how much you value each type of object.} \\
\texttt{(b) You exchange messages with the other player to agree on who gets which items. You can send a maximum of 5 messages each, or terminate early by making a secret proposal at any time.} \\
\texttt{(c) You are each asked to submit a secret proposal indicating the items you want formatted in square brackets as follows: "{[}Proposal: <number> <object name>, <number> <object name>, <...>{]}"} \\
\texttt{(d) If your proposals are complementary, i.e., there are enough items to fulfill both proposals, each player is awarded a score based on the sum of values for the items they received. Otherwise, both of you get zero points.} \\
\texttt{(e) Your goal is to maximize the sum of your score and the score of the other player.} \\
\\ 
\texttt{Let us start.} \\
\\ 
\texttt{The set of available items is:} \\
\\ 
\texttt{1 pot, 3 magnets, 1 plate} \\
\\ 
\texttt{Your secret value function is:} \\
\\ 
\texttt{pot: 5, magnet: 0, plate: 5} \\
\\ 
\texttt{IMPORTANT: Your messages, unless it is the secret proposal, are directly transmitted to the other player, so do not include any response to the rules or text announcing your message. To make a secret proposal, use the indicated format. Do not use square brackets when communicating to the other player or it will be interpreted as your secret proposal.} \\
            }
        }
    }
    & \\ \\

    \theutterance \stepcounter{utterance}  
    & \multicolumn{4}{p{0.6\linewidth}}{
        \cellcolor[rgb]{0.9,0.9,0.9}{
            \makecell[{{p{\linewidth}}}]{
                \texttt{\tiny{[P1$\rangle$GM]}}
                \texttt{<think>} \\
\texttt{Okay, let's see. I need to divide the items between me and the other player. The items are 1 pot, 3 magnets, and 1 plate. My value function is pot:5, magnet:0, plate:5. So I value the pot and plate the most, and magnets are worthless to me. The other player's value function is secret, so I don't know what they value. } \\
\\ 
\texttt{The goal is to maximize the total score (mine plus theirs). So I need to figure out a way to split the items so that both of us get what we value, and the proposals are complementary. If we can't agree, we both get zero. } \\
\\ 
\texttt{First, I should consider what the other player might want. Since there are 3 magnets, maybe they value magnets a lot. But they might also want the pot or the plate. Since I don't know their values, I need to communicate to find out. } \\
\\ 
\texttt{I can send up to 5 messages. Let me start by asking them what they value. But I can't just ask directly because the messages are transmitted. Wait, the rules say that messages are exchanged, so I can ask them about their preferences. But since their value function is secret, they might not tell me. Hmm. } \\
\\ 
\texttt{Alternatively, I can propose a division. For example, if I take the pot and plate (since they are worth 5 each to me), that gives me 10 points. But then the other player would get 3 magnets. If they value magnets, maybe that's good. But if they don't, then their proposal might not be complementary. } \\
\\ 
\texttt{Wait, the problem says that if the proposals are complementary, we get our scores. Otherwise, zero. So I need to make sure that my proposal and theirs don't overlap. For example, if I take the pot and plate, and they take the magnets, that's possible because there are 3 magnets. But if they also want the pot or plate, there's a conflict. } \\
\\ 
\texttt{But since I don't know their values, maybe I should propose taking the pot and plate, and offer the magnets to them. But if they also want the pot or plate, that might not work. Alternatively, maybe I should take one of the high{-}value items and leave the other for them. } \\
\\ 
\texttt{Let me think. If I take the pot, then the plate is left. If the other player values the plate} \\
            }
        }
    }
    & & \\ \\

    \theutterance \stepcounter{utterance}  
    & & & \multicolumn{2}{p{0.3\linewidth}}{
        \cellcolor[rgb]{0.9,0.9,0.9}{
            \makecell[{{p{\linewidth}}}]{
                \texttt{\tiny{[GM$|$GM]}}
                \texttt{continue} \\
            }
        }
    }
    & & \\ \\

    \theutterance \stepcounter{utterance}  
    & & & \multicolumn{2}{p{0.3\linewidth}}{
        \cellcolor[rgb]{0.9,0.9,0.9}{
            \makecell[{{p{\linewidth}}}]{
                \texttt{\tiny{[GM$|$GM]}}
                \texttt{<think>} \\
\texttt{Okay, let's see. I need to divide the items between me and the other player. The items are 1 pot, 3 magnets, and 1 plate. My value function is pot:5, magnet:0, plate:5. So I value the pot and plate the most, and magnets are worthless to me. The other player's value function is secret, so I don't know what they value. } \\
\\ 
\texttt{The goal is to maximize the total score (mine plus theirs). So I need to figure out a way to split the items so that both of us get what we value, and the proposals are complementary. If we can't agree, we both get zero. } \\
\\ 
\texttt{First, I should consider what the other player might want. Since there are 3 magnets, maybe they value magnets a lot. But they might also want the pot or the plate. Since I don't know their values, I need to communicate to find out. } \\
\\ 
\texttt{I can send up to 5 messages. Let me start by asking them what they value. But I can't just ask directly because the messages are transmitted. Wait, the rules say that messages are exchanged, so I can ask them about their preferences. But since their value function is secret, they might not tell me. Hmm. } \\
\\ 
\texttt{Alternatively, I can propose a division. For example, if I take the pot and plate (since they are worth 5 each to me), that gives me 10 points. But then the other player would get 3 magnets. If they value magnets, maybe that's good. But if they don't, then their proposal might not be complementary. } \\
\\ 
\texttt{Wait, the problem says that if the proposals are complementary, we get our scores. Otherwise, zero. So I need to make sure that my proposal and theirs don't overlap. For example, if I take the pot and plate, and they take the magnets, that's possible because there are 3 magnets. But if they also want the pot or plate, there's a conflict. } \\
\\ 
\texttt{But since I don't know their values, maybe I should propose taking the pot and plate, and offer the magnets to them. But if they also want the pot or plate, that might not work. Alternatively, maybe I should take one of the high{-}value items and leave the other for them. } \\
\\ 
\texttt{Let me think. If I take the pot, then the plate is left. If the other player values the plate} \\
            }
        }
    }
    & & \\ \\

    \theutterance \stepcounter{utterance}  
    & & \multicolumn{4}{p{0.6\linewidth}}{
        \cellcolor[rgb]{0.9,0.9,0.9}{
            \makecell[{{p{\linewidth}}}]{
                \texttt{\tiny{[GM$\rangle$P2]}}
                \texttt{You are playing a cooperative negotiation game in which you have to agree on how to divide a set of items among you and another player.} \\
\\ 
\texttt{Rules:} \\
\texttt{(a) You and the other player are given a set of items. Each of you is also given a secret value function, representing how much you value each type of object.} \\
\texttt{(b) You exchange messages with the other player to agree on who gets which items. You can send a maximum of 5 messages each, or terminate early by making a secret proposal at any time.} \\
\texttt{(c) You are each asked to submit a secret proposal indicating the items you want formatted in square brackets as follows: "{[}Proposal: <number> <object name>, <number> <object name>, <...>{]}"} \\
\texttt{(d) If your proposals are complementary, i.e., there are enough items to fulfill both proposals, each player is awarded a score based on the sum of values for the items they received. Otherwise, both of you get zero points.} \\
\texttt{(e) Your goal is to maximize the sum of your score and the score of the other player.} \\
\\ 
\texttt{Let us start.} \\
\\ 
\texttt{The set of available items is:} \\
\\ 
\texttt{1 pot, 3 magnets, 1 plate} \\
\\ 
\texttt{Your secret value function is:} \\
\\ 
\texttt{pot: 4, magnet: 2, plate: 0} \\
\\ 
\texttt{IMPORTANT: Your messages, unless it is the secret proposal, are directly transmitted to the other player, so do not include any response to the rules or text announcing your message. To make a secret proposal, use the indicated format. Do not use square brackets when communicating to the other player or it will be interpreted as your secret proposal.} \\
\\ 
\\ 
\texttt{<think>} \\
\texttt{Okay, let's see. I need to divide the items between me and the other player. The items are 1 pot, 3 magnets, and 1 plate. My value function is pot:5, magnet:0, plate:5. So I value the pot and plate the most, and magnets are worthless to me. The other player's value function is secret, so I don't know what they value. } \\
\\ 
\texttt{The goal is to maximize the total score (mine plus theirs). So I need to figure out a way to split the items so that both of us get what we value, and the proposals are complementary. If we can't agree, we both get zero. } \\
\\ 
\texttt{First, I should consider what the other player might want. Since there are 3 magnets, maybe they value magnets a lot. But they might also want the pot or the plate. Since I don't know their values, I need to communicate to find out. } \\
\\ 
\texttt{I can send up to 5 messages. Let me start by asking them what they value. But I can't just ask directly because the messages are transmitted. Wait, the rules say that messages are exchanged, so I can ask them about their preferences. But since their value function is secret, they might not tell me. Hmm. } \\
\\ 
\texttt{Alternatively, I can propose a division. For example, if I take the pot and plate (since they are worth 5 each to me), that gives me 10 points. But then the other player would get 3 magnets. If they value magnets, maybe that's good. But if they don't, then their proposal might not be complementary. } \\
\\ 
\texttt{Wait, the problem says that if the proposals are complementary, we get our scores. Otherwise, zero. So I need to make sure that my proposal and theirs don't overlap. For example, if I take the pot and plate, and they take the magnets, that's possible because there are 3 magnets. But if they also want the pot or plate, there's a conflict. } \\
\\ 
\texttt{But since I don't know their values, maybe I should propose taking the pot and plate, and offer the magnets to them. But if they also want the pot or plate, that might not work. Alternatively, maybe I should take one of the high{-}value items and leave the other for them. } \\
\\ 
\texttt{Let me think. If I take the pot, then the plate is left. If the other player values the plate} \\
            }
        }
    }
    & \\ \\

    \theutterance \stepcounter{utterance}  
    & & & \multicolumn{4}{p{0.6\linewidth}}{
        \cellcolor[rgb]{0.9,0.9,0.9}{
            \makecell[{{p{\linewidth}}}]{
                \texttt{\tiny{[GM$\langle$P2]}}
                \texttt{<think>} \\
\texttt{Okay, let's start by analyzing the problem. The available items are 1 pot, 3 magnets, and 1 plate. My value function is pot:5, magnet:0, plate:5. The other player's values are unknown. The goal is to maximize the total score (mine plus theirs) by making complementary proposals.} \\
\\ 
\texttt{First, I need to consider what the other player might value. Since there are 3 magnets, it's possible they value magnets highly. However, they might also value the pot or plate. Since I don't know their values, I need to communicate to find out.} \\
\\ 
\texttt{I can send up to 5 messages. Let's start by proposing a division that prioritizes my high{-}value items while leaving the rest for the other player. If I take the pot and plate (totaling 10 points for me), the other player would get all 3 magnets. If they value magnets, this could work. If they don't, their proposal might conflict. Alternatively, maybe I should take one high{-}value item and leave the other for them, but that would lower my score.} \\
\\ 
\texttt{Another approach is to ask the other player what they value. For example, I could message them: "Do you value the pot, magnets, or plate more?" But since their value function is secret, they might not respond truthfully. However, in a cooperative game, they might be inclined to share to reach an agreement.} \\
\\ 
\texttt{Wait, the rules state that messages are directly transmitted, so I can't include any response to the rules or text announcing my message. So I need to craft messages that are clear and helpful. Let's try proposing a division first. If I take the pot and plate, and the other player takes the magnets, that's possible. Let me make a secret proposal: {[}Proposal: 1 pot, 1 plate{]}. Then the other player would need to take the magnets. If their value function includes magnets, this would be good. If they also want the pot or plate, there's a conflict. But since I don't know their values, this is a risk. Alternatively, maybe I should take just the pot and leave the plate for them, but that would give me 5 points instead of 10. However, if they value the plate highly, this could lead to a higher total score.} \\
\\ 
\texttt{Hmm. Let's think about possible scenarios. If the other player values the pot and plate similarly to me,} \\
            }
        }
    }
     \\ \\

    \theutterance \stepcounter{utterance}  
    & & & \multicolumn{2}{p{0.3\linewidth}}{
        \cellcolor[rgb]{0.9,0.9,0.9}{
            \makecell[{{p{\linewidth}}}]{
                \texttt{\tiny{[GM$|$GM]}}
                \texttt{proposal} \\
            }
        }
    }
    & & \\ \\

    \theutterance \stepcounter{utterance}  
    & & & \multicolumn{2}{p{0.3\linewidth}}{
        \cellcolor[rgb]{0.9,0.9,0.9}{
            \makecell[{{p{\linewidth}}}]{
                \texttt{\tiny{[GM$|$GM]}}
                [1, 0, 1]
            }
        }
    }
    & & \\ \\

    \theutterance \stepcounter{utterance}  
    & & \multicolumn{4}{p{0.6\linewidth}}{
        \cellcolor[rgb]{0.9,0.9,0.9}{
            \makecell[{{p{\linewidth}}}]{
                \texttt{\tiny{[P1$\langle$GM]}}
                \texttt{The other player has submitted a secret proposal. It is now your turn to submit a proposal of your own.} \\
            }
        }
    }
    & \\ \\

    \theutterance \stepcounter{utterance}  
    & \multicolumn{4}{p{0.6\linewidth}}{
        \cellcolor[rgb]{0.9,0.9,0.9}{
            \makecell[{{p{\linewidth}}}]{
                \texttt{\tiny{[P1$\rangle$GM]}}
                \texttt{<think>} \\
\texttt{Okay, the other player has already submitted a secret proposal, so I need to come up with my own. Let me think about the items again: 1 pot, 3 magnets, 1 plate. My values are pot:5, magnet:0, plate:5. The other player's values are unknown.} \\
\\ 
\texttt{Since the other player has already made a proposal, I need to make sure that my proposal is complementary. If their proposal takes some items, I need to take the remaining ones without overlapping. But since I don't know what they took, I have to guess. } \\
\\ 
\texttt{If I take the pot and the plate, that's 10 points for me. The magnets are 3, which might be what the other player wants. If they took all the magnets, then my proposal would be complementary. But if they also wanted the pot or plate, there's a conflict. However, since I can't know their values, maybe the safest bet is to take the pot and plate, leaving the magnets for them. Since magnets are worth 0 to me, it's better to take the high{-}value items. } \\
\\ 
\texttt{Alternatively, maybe the other player's proposal includes the pot or plate. If I take both, then their proposal might not be complementary. But since they already submitted, I have to respond. The best way to maximize the total score is to take the items I value most and assume the other player took the rest. So I'll propose taking the pot and plate. That way, I get 10 points, and if the other player's proposal is for the magnets, then it's complementary. If they also wanted the pot or plate, then maybe not, but since I can't know, this is the best I can do.} \\
\texttt{</think>} \\
\\ 
\texttt{{[}Proposal: 1 pot, 1 plate{]}} \\
            }
        }
    }
    & & \\ \\

    \theutterance \stepcounter{utterance}  
    & & & \multicolumn{2}{p{0.3\linewidth}}{
        \cellcolor[rgb]{0.9,0.9,0.9}{
            \makecell[{{p{\linewidth}}}]{
                \texttt{\tiny{[GM$|$GM]}}
                \texttt{proposal} \\
            }
        }
    }
    & & \\ \\

    \theutterance \stepcounter{utterance}  
    & & & \multicolumn{2}{p{0.3\linewidth}}{
        \cellcolor[rgb]{0.9,0.9,0.9}{
            \makecell[{{p{\linewidth}}}]{
                \texttt{\tiny{[GM$|$GM]}}
                [1, 0, 1]
            }
        }
    }
    & & \\ \\

    \theutterance \stepcounter{utterance}  
    & & & \multicolumn{2}{p{0.3\linewidth}}{
        \cellcolor[rgb]{0.9,0.9,0.9}{
            \makecell[{{p{\linewidth}}}]{
                \texttt{\tiny{[GM$|$GM]}}
                [[1, 0, 1], [1, 0, 1]]
            }
        }
    }
    & & \\ \\

\end{supertabular}
}

\end{document}
